\section{Local gates from boundary values rather than derivative queries}
\label{sec:value-probe}
The effective fingerprint permits a weaker local measurement primitive.
We keep the launch and trace experiment, the local nature of all queries,
and the distance-to-solid test. We no longer require measured derivatives
or an intrinsic arclength oracle. Neither this reduction nor its proof
turns active geometric access into a passive trajectory observation.

\subsection{The value-access model}
At an encountered impact the apparatus supplies a Cartesian graph chart
on a certified interval of fixed positive length, with a smaller working
interval and known transverse-coordinate direction. Chart origins and
axes are laboratory coordinates, not unknown type or period labels.
At any requested abscissa in that chart it returns a height enclosure of
width at most $2\epsilon$. An error in the requested abscissa is charged
as height error using the known slope bound. The working chart and a
collar have a numerical $C^3$ bound $M_3$, enlarged to bound all lower
derivatives. The uniform analytic and curvature priors supply such
bounds after the chart radius is reduced. The chart domains, coverage
and nonverticality margins are part of this local access model.

The apparatus also encloses distance to the solid at requested points
in the permitted neighborhood of the candidate flight, with stated
absolute error. This query tests all solid in that neighborhood; it is
not a query only against the two already encountered patches. It remains
a genuine additional measurement. No complete boundary, body catalogue,
period, translated-copy label, derivative, or arclength value is returned.
Let $V(\epsilon)\ge1$ be a calibrated upper bound for the cost of a
height or distance enclosure of accuracy $\epsilon$, nonincreasing as a function of $\epsilon$, so its cost rises as
accuracy improves. Bit complexity of the subsequent computation is
separate unless explicitly included in that bound.

\begin{lemma}[Certified differentiation of a local value chart]
\label{lem:value-differences}
Suppose $f$ has third derivative bounded by $M_3$ on $[x-t,x+t]$,
and the three height measurements there have error at most $\epsilon$.
Their centered first and second differences $D_1,D_2$ satisfy
\begin{align}
 |D_1-f'(x)|&\le M_3t^2/6+\epsilon/t,\label{eq:value-first}\\
 |D_2-f''(x)|&\le M_3t/3+4\epsilon/t^2.\label{eq:value-second}
\end{align}
On a fixed smaller chart interval, finitely many height values produce
a smooth approximation and certified uniform errors through order two
of at most $C\nu$. It suffices to use $O(1+\nu^{-1})$ values of
accuracy $c\nu^3$, where $C,c>0$ depend only on the chart bounds.
\end{lemma}
\begin{proof}
Taylor's formula with integral remainder bounds each third-order
remainder by $M_3t^3/6$. Subtract the two Taylor formulas for the first
difference and add them for the second. The coefficient absolute sums
of the measurement errors are $1/t$ and $4/t^2$, respectively. This
proves the displayed bounds without assuming any stochastic noise model.

Choose a uniform chart grid of spacing $t=c_0\nu$ inside the collar.
Fit the quadratic polynomial through three neighboring height values.
Its degree-two Taylor comparison at the middle point has derivative
error at most $C M_3t^{3-j}+C\epsilon t^{-j}$ for $0\le j\le2$
on a fixed enlargement of that cell. Patch these polynomials with a
smooth partition of unity of bounded overlap, whose $j$th derivatives
are $O(t^{-j})$. The product rule gives the same errors for the patched
approximation. Taking $\epsilon\le c_1\nu^3$ proves the claim.
The constants may be bounded from the three-node interpolation matrix,
a fixed chosen partition, the interval length and $M_3$. Alternatively,
\eqref{eq:value-first}--\eqref{eq:value-second} with
$t\le\min(1,\nu/[4(1+M_3)])$ and
$\epsilon\le\nu t^2/32$ give pointwise enclosures; the prior third
and second derivative bounds extend them between grid points.
\end{proof}

\begin{theorem}[Value-probed construction of periodic local records]
\label{thm:value-gates}
On the uniform class, with the numerical priors in
Section~\ref{sec:effective}, replace the derivative-and-arclength probe
by the value-access model above. A finite local procedure has the
validity and protected-return properties of Lemma~\ref{lem:probe}.
It produces $S_N$ to error less than $\chi_*/16$ and intrinsic endpoint
coordinates to any prescribed error $\nu>0$. The registry, physical
hazard, finite-list certificate and stopping conclusions remain valid,
with constants depending on the same prior and the chart-access margins.

After reducing $\nu$ below fixed prior-dependent thresholds, the number
of height or distance queries is at most $C(1+N)(1+\nu^{-1})$, each
requiring accuracy no smaller than $c\nu^3$. In particular their total
measurement cost is bounded by
\begin{equation}\label{eq:value-cost}
 P_{\rm val}(\nu)\le C(1+N)(1+\nu^{-1})V(c\nu^3).
\end{equation}
The large but fixed $N$ is given by \eqref{eq:eff-integers}; it is
not a derivative order or a growing histogram dimension. This is an
upper bound on local measurements, not a polynomial bound in the prior
margins or in the bit complexity of analytic fitting.
\end{theorem}
\begin{proof}
\emph{Local normal certification.}
Apply Lemma~\ref{lem:value-differences} on the slightly larger
encountered charts. The distance function between two chart points and
its first two derivatives are continuous rational expressions in the
chart values, their first two derivatives and the positive distance.
The gap and chart-transversality margins therefore give certified
$O(\nu)$ enclosures for these expressions. The normal Hessian in
\eqref{eq:normal-hessian}, conjugated by the bounded changes from graph
coordinates to arclength, is uniformly positive in a common neighborhood.
Its modulus of continuity is bounded using the $C^3$ prior.

One explicit existence test uses a fixed approximate inverse $B_1$ of
the Hessian and the map $z\mapsto z-B_1\nabla D(z)$. On a small
closed ball, interval bounds certify both a contraction factor less than
one and that the ball maps into itself. Banach's theorem then certifies
a unique normal point there. Interval subdivision and refinement of
chart values make this test pass on the protected neighborhood: its true
normal point is interior, the Hessian has a fixed positive lower bound,
and the inverse and continuity bounds are uniform. Refinement of the
same contraction gives a normal-point enclosure of size $O(\nu)$.
Tests without certified slack return no record under a fixed finite
budget. A fixed number of coarse balls covers the allowed working
region; the refinement depth for a protected input is logarithmic in
$1/\nu$ and is dominated by the stated grid budget. No derivative
returned by a sensor is used in this test.

\emph{Clearance and canonical coordinates.}
The supporting-half-plane proof of the global shortest property still
applies to the certified facing normal. Use the same finite
1-Lipschitz distance-to-solid mesh as in Section~\ref{sec:acquisition},
with strict error and clearance slack. Its number of queries and required
accuracy are fixed by the prior. Endpoint collars are controlled by the
local chart enclosures, convexity and the separation bound. Reducing
$\nu$ makes the normal root, gap and canonical frame errors at most
$C\nu$. On the fixed smaller canonical graph interval, nonverticality
allows conversion of a value chart to the normal graph by one-dimensional
monotone inversion. Uniform derivative bounds give condition numbers
bounded by prior constants. The canonical graph values obtained from
the approximating charts differ from their true values by at most
$C\nu$, uniformly on the whole interval.

Consequently evaluation at the fixed $N+1$ nodes on each end gives a
fingerprint error at most $C\nu$. Reduce the working accuracy below
$\chi_*/(32C)$, including the small error from the approximating gap.
Theorem~\ref{thm:effective} and Proposition~\ref{prop:registry} give
exact class and relative sign recognition. The grid used for graph
approximation need not resolve the spacing between every fingerprint
node: its certified uniform approximation can be evaluated at every
node. Counting a separate refinement per node would still satisfy the
more conservative factor $1+N$ in \eqref{eq:value-cost}.

\emph{Intrinsic offsets and physical returns.}
The arclength integrand of a graph is $\sqrt{1+(f')^2}$. It is
1-Lipschitz as a function of $f'$ and has a bounded spatial derivative
under the $C^2$ prior. Thus the derivative enclosures and a Riemann mesh
of spacing $c\nu$ give its integral, from the enclosed normal point
to each measured endpoint, with certified error $C\nu$. The uncertainty
of that normal point and of the endpoint position contributes the same
order. Rescale the internal tolerance to obtain the requested final
error $\nu$. There is no supplied exact intrinsic origin or arclength
measurement in this step.

These constructions have precisely the margins used in the local
validity test. On its protected return rectangle every branch is
recognized, and no foreign branch is accepted; ambiguous inputs are
rejected. Choose the histogram rectangle and a collar strictly inside
that protected rectangle. Borderline measured endpoints can affect
only cell-boundary bands, including the outer histogram boundary.
The bound in Section~\ref{sec:statistics} charges those errors. The
ideal recognition rule is translation invariant because the canonical
values and intrinsic offsets are, while bounded nonperiodic numerical
errors do not change the registry decision. Fresh launches therefore
have the same positive protected-return hazard, up to reducing the
fixed constants. The inherited inverse and certificate apply unchanged.

Finally the chart grids use $O(1+\nu^{-1})$ values at accuracy
$c\nu^3$. A fixed number of charts and clearance meshes suffices for
a bounded-time trace. Root refinement, chart conversion and quadrature
can all use these certified approximations; allowing additional value
queries per fingerprint node gives the displayed conservative count.
This proves \eqref{eq:value-cost}. Arithmetic on the approximations
and evaluating many fingerprint coordinates are real additional work;
the theorem's measurement bound does not suppress them in a bit-cost claim.
\end{proof}

\subsection{Tail exponents and the cost of the weaker probe}
The stopping schedule uses a chosen $a\ge6$ throughout. If a local
probe has cost $P(\nu)\le C\nu^{-s}$, epoch $k$ requires at most
$Ck^{3+s}$ such work. Its execution is decided by the past, so summing
this bound against the tail of $T_\xi$ at $k-1$ gives a finite expectation
whenever $a>4+s$. At the fixed choice $a=6$, this argument only covers
$s<2$; for $s=2$ its bound is a harmonic series. This is a limitation
of that cost guarantee, not a proof that every such apparatus has infinite
expected cost.

For value access with $V(\epsilon)\le C\epsilon^{-s_v}$,
\eqref{eq:value-cost} gives $s=1+3s_v$. Thus through epoch $m$ the
measurement cost is $O(m^{5+3s_v})$, and the sufficient condition is
\begin{equation}\label{eq:value-expectation}
                       a\ge6,\qquad a>5+3s_v.
\end{equation}
The launch count, laboratory range and statistical powers do not change.
Larger $a$ changes the logarithmic error-spending factors, not those
powers. Apparatus movement, arbitrary unbounded query costs and fitting
complexity still require their own accounts. The original stronger-probe
result remains a separate valid specialization of the same schedule.
